\begin{enumerate}
\def\labelenumi{\arabic{enumi}.}
\setcounter{enumi}{3}
\item
  \textbf{Detailed Formatting Instructions}
\end{enumerate}

The styles and formats for the AIAA Papers Template have been
incorporated into the structure of this document. If you are using
Microsoft Word 2001 or later, please use this template to prepare your
manuscript. Regardless of which program you use to prepare your
manuscript, please use the formatting instructions contained in this
document as a guide.

\subsection{Document Text}\label{document-text}

The default font for the Project Technical Report is Times New Roman,
10-point size. In the electronic template, use the ``Text'' or
``Normal'' style from the pull-down menu to format all primary text for
your manuscript. The first line of every paragraph should be indented,
and all lines should be single-spaced. Default margins are 1'' on all
sides. In the electronic version of this template, all margins and other
formatting is preset. There should be no additional lines between
paragraphs.

\begin{quote}
Extended quotes, such as this example, are to be used when material
being cited is longer than a few sentences, or the standard quotation
format is not practical. In this Word template, the appropriate style is
``Extended Quote'' from the drop-down menu. Extended quotes are to be in
Times New Roman, 9-point font, indented 0.4'' and full justified.
\end{quote}

\emph{NOTE:} If you are using the electronic template to format your
manuscript, the required spacing and formatting will be applied
automatically, simply by using the appropriate style designation from
the pull-down menu.

\subsection{Headings}\label{headings}

The title of your paper should be typed in bold, 24-point type, with
capital and lower-case letters, and centered at the top of the page. The
names of the authors, business or academic affiliation, city, and
state/province should follow on separate lines below the title. The
names of authors with the same affiliation can be listed on the same
line above their collective affiliation information. Author names are
centered, and affiliations are centered and in italic type immediately
below the author names. The affiliation line for each author is to
include that author's city, state, and zip/postal code (or city,
province, zip/postal code and country, as appropriate). The first-page
footnotes (lower left-hand side) contain the job title and department
name, and street address/mail stop for each author. Author email
addresses may be included also.

Major headings (``Heading 1'' in the template style list) are bold
11-point font, flush-left, and numbered with numerals.

\begin{quote}
Subheadings (``Heading 2'' in the template style list) are bold, flush
left, and numbered with capital letters. Sub-Subheadings (``Heading 3''
in the template style list) are italic, flush left, and numbered (1. 2.
3. etc.)
\end{quote}

\subsection{Abstract}\label{abstract}

The abstract should appear at the beginning of your paper. It should be
one paragraph long (not an introduction) and complete (no reference
numbers). It should indicate subjects dealt with in the paper and state
the objectives of the investigation. Newly observed facts and
conclusions of the experiment or argument discussed in the paper must be
stated in summary form; readers should not have to read the paper to
understand the abstract. The abstract should be bold, indented 3 picas
(1/2'') on each side, and separated from the rest of the document by -
blank lines above and below the abstract text.

\subsection{Nomenclature}\label{nomenclature}

Papers with many symbols may benefit from a nomenclature list that
defines all symbols with units, inserted between the abstract and the
introduction. If one is used, it must contain all the symbology used in
the manuscript, and the definitions should not be repeated in the text.
In all cases, identify the symbols used if they are not widely
recognized in the profession. Define acronyms in the text, not in the
nomenclature.

\subsection{Footnotes and References}\label{footnotes-and-references}

Footnotes, where they appear, should be placed above the 1'' margin at
the bottom of the page. To insert footnotes into the template, use the
Insert\textgreater Footnote feature from the main menu as necessary.
Numbered footnotes as formatted automatically in the template are
acceptable, but superscript symbols are the preferred AIAA style, in the
sequence, *, †, ‡, §, ¶, \#, **. ††, ‡‡, §§, etc.

List and number all references at the end of the paper. Corresponding
bracketed numbers are used to cite references in the text {[}1{]},
unless the citation is an integral part of the sentence (e.g., ``It is
shown in Ref. {[}2{]} that\ldots'') or follows a mathematical
expression: ``A\textsuperscript{2} + B = C (Ref. {[}3{]}).'' For
multiple citations, separate reference numbers with commas {[}4, 5{]},
or use a dash to show a range {[}6-8{]}. Reference citations in the text
should be in numerical order.

In the reference list, give all authors' names; do not use ``et
al\emph{.}'' unless there are more than 10 authors. Papers that have not
been published should be cited as ``unpublished''; papers that have been
submitted or accepted for publication should be cited as ``submitted for
publication.'' Private communications and personal websites should
appear as footnotes rather than in the reference list.

References should be cited according to the standard publication
reference style (for examples, see the ``References'' section of this
template). Never edit titles in references to conform to AIAA style of
spellings, abbreviations, etc. Names and locations of publishers should
be listed; month and year should be included for reports and papers. For
papers published in translation journals, please give the English
citation first, followed by the original foreign language citation.

\subsection{Images, Figures, and
Tables}\label{images-figures-and-tables}

All artwork, captions, figures, graphs, and tables will be reproduced
exactly as submitted. Be sure to position any figures, tables, graphs,
or pictures as you want them printed.

Do not insert your tables and figures in text boxes. Figures should have
no background, borders, or outlines. In the electronic template, use the
``Figure'' style from the pull-down formatting menu to type caption
text. You may also insert the caption by going to the References menu
and choosing Insert Caption. Make sure the label is ``Fig.,'' and type
your caption text in the box provided. Captions are bold with a single
tab (no hyphen or other character) between the figure number and figure
description.

\includegraphics[width=3.45833in,height=2.63194in,alt={Image}]{./media/image1.jpg}

\textbf{Fig. 1. Magnetization as a function of applied fields.}

Place figure captions below all figures; place table titles above the
tables. If your figure has multiple parts, include the labels ``a),''
``b),'' etc. below and to the left of each part, above the figure
caption. Please verify that the figures and tables you mention in the
text exist. \emph{Please do not include captions as part of the figures,
and do not put captions in separate text boxes linked to the figures}.
When citing a figure in the text, use the abbreviation ``Fig.'' except
at the beginning of a sentence. Do not abbreviate ``Table.'' Number each
different type of illustration (i.e., figures, tables, images)
sequentially with relation to other illustrations of the same type.

Figure axis labels are often a source of confusion. Use words rather
than symbols. As in the example to the right, write the quantity
``Magnetization'' rather than just ``M.'' Do not enclose units in
parenthesis, but rather separate them from the preceding text by commas.
Do not label axes only with units. As in Fig. 1, for example, write
``Magnetization, kA/m'' not just ``kA/m.'' Do not label axes with a
ratio of quantities and units. For example, write ``Temperature, K,''
not ``Temperature/K.''

Multipliers can be especially confusing. Write ``Magnetization, kA/m''
or ``Magnetization, 10\textsuperscript{3} A/m.'' Do not write
``Magnetization (A/m) x 1000'' because the reader would not then know
whether the top axis label in Fig. 1 meant 16000 A/m or 0.016 A/m.
Figure labels must be legible, and all text within figures should be
uniform in style and size, no smaller than 8-point type.

\subsection{Equations, Numbers, Symbols, and
Abbreviations}\label{equations-numbers-symbols-and-abbreviations}

Equations are centered and numbered consecutively, with equation numbers
in parentheses flush right, as in Eq. (1). Insert a blank line above and
below the equation. First use the equation editor to create the
equation. If you are using Microsoft Word, use either the Microsoft
Equation Editor or the MathType add-on
(\href{http://www.mathtype.com}{\ul{http://www.mathtype.com}}) for
equations in your paper, use the function
(Insert\textgreater Object\textgreater Create New\textgreater Microsoft
Equation \emph{or} MathType Equation) to insert it into the document.
Please note that ``Float over text'' should \emph{not} be selected. To
insert the equation into the document:

\begin{enumerate}
\def\labelenumi{\arabic{enumi})}
\item
  Select the ``Equation'' style from the pull-down formatting menu and
  hit ``tab'' once.
\item
  Insert the equation, hit ``tab'' again,
\item
  Enter the equation number in parentheses.
\end{enumerate}

A sample equation is included here, formatted using the preceding
instructions. To make your equation more compact, you can use the
solidus (/), the exp function, or appropriate exponents. Use parentheses
to avoid ambiguities in denominators.

\includegraphics[width=3.09028in,height=0.54861in,alt={Image2}]{./media/image2.jpg}
(1)

Be sure that the symbols in your equation are defined before the
equation appears or immediately following. Italicize symbols (\emph{T}
might refer to temperature, but T is the unit tesla). Refer to ``Eq.
(1),'' not ``(1)'' or ``equation (1)'' except at the beginning of a
sentence: ``Equation (1) is\ldots'' Equations can be labeled other than
``Eq.'' should they represent inequalities, matrices, or boundary
conditions. If what is represented is more than one equation, the
abbreviation ``Eqs.'' can be used.

Define abbreviations and acronyms the first time they are used in the
text, even after they have already been defined in the abstract. Very
common abbreviations such as AIAA, SI, ac, and dc do not have to be
defined. Abbreviations that incorporate periods should not have spaces:
write ``P.R.,'' not ``P. R.'' Delete periods between initials if the
abbreviation has three or more initials; e.g., U.N. but ESA. Do not use
abbreviations in the title unless they are unavoidable (for instance,
``AIAA'' in the title of this article).

\subsection{General Grammar and Preferred
Usage}\label{general-grammar-and-preferred-usage}

Use only one space after periods or colons. Hyphenate complex modifiers:
``zero-field-cooled magnetization.'' Avoid dangling participles, such
as, ``Using Eq. (1), the potential was calculated.'' {[}It is not clear
who or what used Eq. (1).{]} Write instead ``The potential was
calculated using Eq. (1),'' or ``Using Eq. (1), we calculated the
potential.''

Use a zero before decimal points: ``0.25,'' not ``.25.'' Use
``cm\textsuperscript{2},'' not ``cc.'' Indicate sample dimensions as
``0.1 cm x 0.2 cm,'' not ``0.1 x 0.2 cm\textsuperscript{2}.'' The
preferred abbreviation for ``seconds'' is ``s,'' not ``sec.'' Do not mix
complete spellings and abbreviations of units: use
``Wb/m\textsuperscript{2}'' or ``webers per square meter,'' not
``webers/m\textsuperscript{2}.'' When expressing a range of values,
write ``7 to 9'' or ``7-9,'' not ``7\textasciitilde9.''

A parenthetical statement at the end of a sentence is punctuated outside
of the closing parenthesis (like this). (A parenthetical sentence is
punctuated within parenthesis.) In American English, periods and commas
are placed within quotation marks, like ``this period.'' Other
punctuation is ``outside''! Avoid contractions; for example, write ``do
not'' instead of ``don't.'' The serial comma is preferred: ``A, B, and
C'' instead of ``A, B and C.''

If you wish, you may write in the first person singular or plural and
use the active voice (``I observed that\ldots'' or ``We observed
that\ldots'' instead of ``It was observed that\ldots''). Remember to
check spelling. If your native language is not English, please ask a
native English-speaking colleague to proofread your paper.

The word ``data'' is plural, not singular (i.e., ``data are,'' not
``data is''). The subscript for the permeability of vacuum
µ\textsubscript{0} is zero, not a lowercase letter ``o.'' The term for
residual magnetization is ``remanence''; the adjective is ``remanent'';
do not write ``remnance'' or ``remnant.'' The word ``micrometer'' is
preferred over ``micron'' when spelling out this unit of measure. A
graph within a graph is an ``inset,'' not an ``insert.'' The word
``alternatively'' is preferred to the word ``alternately'' (unless you
really mean something that alternates). Use the word ``whereas'' instead
of ``while'' (unless you are referring to simultaneous events). Do not
use the word ``essentially'' to mean ``approximately'' or
``effectively.'' Do not use the word ``issue'' as a euphemism for
``problem.'' When compositions are not specified, separate chemical
symbols by en-dashes; for example, ``NiMn'' indicates the intermetallic
compound Ni\textsubscript{0.5}Mn\textsubscript{0.5} whereas ``Ni--Mn''
indicates an alloy of some composition
Ni\textsubscript{x}Mn\textsubscript{1-x}.

Be aware of the different meanings of the homophones ``affect'' (usually
a verb) and ``effect'' (usually a noun), ``complement'' and
``compliment,'' ``discreet'' and ``discrete,'' ``principal'' (e.g.,
``principal investigator'') and ``principle'' (e.g., ``principle of
measurement''). Do not confuse ``imply'' and ``infer.''

Prefixes such as ``non,'' ``sub,'' ``micro,'' ``multi,'' and ``"ultra''
are not independent words; they should be joined to the words they
modify, usually without a hyphen. There is no period after the ``et'' in
the abbreviation ``et al\emph{.}'' The abbreviation ``i.e.,'' means
``that is,'' and the abbreviation ``e.g.,'' means ``for example'' (these
abbreviations are not italicized).
